% Сборка из EDR_code/ (два прохода):
% xelatex -halt-on-error -interaction=nonstopmode \
%   -output-directory=SVD/synthesis_2026_10_24 SVD/AI_research_synthesis.tex
\documentclass[11pt,a4paper]{article}
\usepackage{fontspec}
\setmainfont{Noto Serif}
\setsansfont{Noto Sans}
\setmonofont[Scale=0.84]{DejaVu Sans Mono}
\usepackage[russian]{babel}
\usepackage[margin=24mm]{geometry}
\usepackage{amsmath,amssymb,amsthm,mathtools}
\usepackage{booktabs,tabularx,longtable,array}
\usepackage{enumitem,microtype,xurl}
\usepackage[unicode,hidelinks]{hyperref}
\usepackage{bookmark}
\hypersetup{pdftitle={Идеи развития ADP и результаты их проверки},
  pdfsubject={Систематизация исследовательских материалов SVD/},pdfauthor={}}
\setlength{\emergencystretch}{3em}
\setlength{\parskip}{3pt}
\setlist{itemsep=2pt,topsep=4pt}
\numberwithin{equation}{section}
\newcommand{\R}{\mathbb R}
\newcommand{\Id}{\mathrm{Id}}
\newcommand{\T}{^{\mathsf T}}
\newcommand{\F}{_{\mathrm F}}
\newcommand{\ip}[2]{\langle #1,#2\rangle_{\mathrm F}}
\newcommand{\file}[1]{\nolinkurl{#1}}
\newcommand{\repo}[1]{\href{../#1}{\nolinkurl{#1}}}
\DeclareMathOperator{\rank}{rank}
\DeclareMathOperator{\grad}{grad}
\DeclareMathOperator{\polar}{polar}
\DeclareMathOperator{\diag}{diag}
\DeclareMathOperator{\spanop}{span}
\newtheorem{proposition}{Утверждение}[section]
\title{Идеи развития ADP и результаты их проверки\\[5pt]
  \large Малоранговые поправки, профильная оптимизация,\\
  геометрия подпространств и адаптивная локализация}
\author{}
\date{Систематизация материалов \texttt{SVD/} по состоянию на 24 октября 2026 года}

\begin{document}
\maketitle
\begin{abstract}
Систематизированы исследовательские записки и обсуждения, подготовленные
с участием языковых моделей и собранные в каталоге \file{SVD/}.
Рассмотрены условный SVD-решатель, малоранговые поправки, профилирование
локальных коэффициентов, шаги на многообразии Грассмана, дополнения Шура,
метрические и тензорные варианты, а также отбор центров локализации.
Математические тождества сопоставлены с численными проверками формул,
измерениями отдельных операций и парными экспериментами полного обучения.
Для относительного шага Гаусса~--- Ньютона имеется независимая положительная
проверка в ограниченном классе двухиндексных задач. Полный касательный шаг,
совместная оптимизация углов и проекционный якобиан не прошли соответствующие
критерии отбора; LoRA-подобная поправка проверена только в пилоте.
Для части чисел, приведённых в исходных рукописях, первичные результаты
не найдены. Эти числа отделены от выводов, подтверждённых сохранёнными
экспериментами. Общая гарантия восстановления EDR-подпространства
и сходимости внешней адаптивной процедуры в рассмотренных материалах
не установлена.
\end{abstract}
\tableofcontents
\newpage

\section{Предмет систематизации и состав свидетельств}

В исследовательском корпусе изучается уточнение малоразмерного
подпространства в процедуре средних производных (ADP) по локальным
линейным статистикам. Обсуждаются свободная матрица и её поправка,
ортонормированный базис и пространственно меняющееся поле подпространств.
Использование SVD во всех этих задачах не делает их эквивалентными.

Исходные обсуждения находятся в \file{archive/}; более поздние тексты
уточняют постановки и исправляют некоторые ранние утверждения.
\file{SVD_solver.tex} систематизирует условную матричную задачу,
а \file{EDR_unified_theory.tex} связывает её с профилированием
и геометрией Грассмана. Записки \file{kimi.tex}, \file{glm.tex},
\file{edr_research_note.tex}, \file{Tucker.md}, \file{argon.md}
и \file{cladue.md} содержат дальнейшие предложения.
\file{multiindex.tex} служит описанием исходной ADP-процедуры.
Текстовые копии в \file{txt/} и PDF-версии не считаются самостоятельными
проверками. Карта происхождения идей приведена в приложении~\ref{sec:corpus}.

В изложении различаются четыре вида результатов.
\emph{Математическое утверждение} имеет сформулированные условия
и вывод из заданного функционала. \emph{Численная проверка формулы}
сопоставляет её с независимым плотным расчётом или конечными разностями
на небольших задачах. \emph{Измерение операции} относится к выбранной
статистике и отдельному вычислительному участку. \emph{Проверка полного
обучения} включает построение статистик и адаптивный внешний цикл,
а качество измеряет относительно известного истинного подпространства.
Последний вид проверки также требует различать предварительный пилот,
отбор на заранее зафиксированных данных и независимую отложенную проверку.

Численные выводы ниже относятся к сохранённым конфигурациям и версиям
кода. Время отдельного ядра, время \file{model.fit} и полное время процесса
с импортами различаются. Оценки памяти через \file{tracemalloc} не равны
пиковому RSS процесса. При подготовке настоящего документа длительные
эксперименты не повторялись; заново проверены малые алгебраические примеры
и агрегаты избранных сохранённых прогонов.

\section{Две внутренние задачи и критерий восстановления}
\label{sec:problems}

\subsection{Обозначения и условная квадратичная задача}

Размерность пространства признаков обозначим $d$, размерность искомого
подпространства~--- $m<d$. Число центров равно $J$, число направленных
моментов на центр~--- $p$. На одной внешней итерации заданы
\[
 U_j\in\R^{p\times d},\quad I_j\in\R^p,\quad N_j>0,
 \qquad j=1,\ldots,J.
\]
Эти величины фиксируются на время внутреннего решения. Матрица
$Y_0\in\R^{d\times m}$ имеет ортонормированные столбцы. В исходниках
SVD часто используется строковый базис $P=Y_0\T$; здесь все базисы
записаны столбцами.

При фиксированных локальных коэффициентах $g_j\in\R^m$ рассматривается
\begin{equation}\label{eq:fixed}
 F_g(B)=\sum_jN_j\|I_j-U_jBg_j\|^2
             +\lambda\|B-Y_0\|\F^2,
 \qquad B\in\R^{d\times m}.
\end{equation}
Ограничения $\rank(B)\le r$, $\rank(B-Y_0)\le r$ и отсутствие
ограничения ранга определяют три разные задачи. При $\lambda>0$
неограниченная задача строго выпукла. Условный минимум по матрице
не является минимумом всего адаптивного ADP: после изменения весов
или коэффициентов меняется сама внутренняя задача.\cite{svd,algorithm}

\subsection{Профильная задача на подпространствах}

Другая постановка исключает коэффициенты заново при каждом базисе:
\begin{equation}\label{eq:profile}
 \Phi_\tau(Y)=\sum_jN_j\min_{g\in\R^m}
 \bigl\{\|I_j-U_jYg\|^2+\tau\|g\|^2\bigr\},
 \qquad Y\T Y=\Id_m.
\end{equation}
Замены $Y\mapsto YQ$, $g\mapsto Q\T g$, $Q\T Q=\Id_m$,
сохраняют функционал: оптимизируется подпространство, то есть точка
многообразия Грассмана $\mathrm{Gr}(m,d)$.
При $\tau>0$ локальные системы невырождены. При $\tau=0$
профиль требует анализа ранга; на границах ранга гладкость может нарушаться.

Параметры $\lambda$ в \eqref{eq:fixed}, $\tau$ в \eqref{eq:profile}
и демпфирование линейной модели шага имеют разный смысл.
Первые два входят в целевой функционал. Демпфирование регулирует
предложенный шаг и само по себе не добавляет статистического штрафа.
В проверенных Grassmann-сравнениях основная цель~--- \eqref{eq:profile}
при $\tau=0$; штраф по расстоянию до $Y_0$ из некоторых рукописей
в эту цель не входит.\cite{unified,relative,glmcheck}

\subsection{Ранг и ошибка подпространства}

Малый ранг самой матрицы $B$ ограничивает число представленных
направлений. Малый ранг поправки $B-Y_0$ ограничивает изменение
базиса. Ранг касательного шага определяет число одновременно
вращаемых плоскостей. Ни один из этих параметров не служит автоматически
оценкой структурной размерности $m$.

В воспроизводимых поздних экспериментах EDR-ошибка определяется как
\begin{equation}\label{eq:edr}
 e(\widehat Y,Y_*)=
 \frac{\|\widehat Y\widehat Y\T-Y_*Y_*\T\|\F}{\sqrt{2m}}
 =\left(\frac1m\sum_{a=1}^m\sin^2\theta_a\right)^{1/2}.
\end{equation}
Здесь $\theta_a$~--- главные углы, оба базиса ортонормированы
и имеют $m$ столбцов. Поэтому уменьшение $\Phi_\tau$ и уменьшение
$e$~--- разные результаты. Внешняя процедура изменяет локализацию
и моменты; её последовательность не минимизирует один неизменный
функционал. Остановка по минимальной ширине окна $h_{\min}$
означает завершение расписания, а не сертификат стационарности.

Для парных сравнений поздних Grassmann-исследований приводятся
геометрическое среднее отношений EDR-ошибок
\[
 R_e=\exp\left(\frac1L\sum_{\ell=1}^L
         \log\frac{e_{\mathrm{cand},\ell}}{e_{\mathrm{base},\ell}}\right)
\]
и медиана отношений времени. Это не отношение отдельных медиан.
Интервалы неопределённости относятся к процедуре ресэмплирования
конкретного исследования; во всех таблицах ниже сохранена исходная
единица парности.\cite{metricdef,relative,combined,lowrank}

\section{Условный SVD-решатель: обоснованные конструкции и ограничения}
\label{sec:fixedsvd}

\subsection{Операторная форма и точный шаг ранга один}

Введём $(\mathcal AX)_j=\sqrt{N_j}U_jXg_j$,
$b_j=\sqrt{N_j}I_j$. Тогда
\begin{align}
 \mathcal A^*z&=\sum_j\sqrt{N_j}U_j\T z_jg_j\T,\nonumber\\
 \mathcal LX&=\sum_jN_jU_j\T U_jXg_jg_j\T+\lambda X,
 &C&=\mathcal A^*b+\lambda Y_0.\label{eq:operator}
\end{align}
При $\lambda>0$ оператор $\mathcal L$ положительно определён.
Полный минимум удовлетворяет $\mathcal LB_*=C$, а
\begin{equation}\label{eq:energy}
 F_g(B)-F_g(B_*)=\ip{B-B_*}{\mathcal L(B-B_*)}.
\end{equation}
Действия через $U_j$ и $U_j\T$ позволяют не строить плотную
матрицу размера $dm\times dm$.

Для выбранного атома $E=va\T$, $\|v\|=\|a\|=1$,
и $Q=C-\mathcal LB$ имеем
\begin{equation}\label{eq:atom}
 t_* =\frac{\ip{Q}{E}}{\ip{E}{\mathcal LE}},\qquad
 F_g(B)-F_g(B+t_*E)=\frac{\ip{Q}{E}^2}{\ip{E}{\mathcal LE}}.
\end{equation}
Это точный выбор масштаба данного атома. SVD невязки $Q$
максимизирует числитель при евклидовой нормировке, но знаменатель
зависит от направлений. Поэтому ведущая сингулярная пара не обязана
давать максимальное фактическое уменьшение.

Исходная жадная схема попеременно уточняет $a$ и $v$, добавляет
компоненту, сжимает накопленные факторы и совместно переоценивает
масштабы. Последняя операция решает малую МНК-задачу в оболочке
выбранных атомов и учитывает их перекрёстные взаимодействия.
QR факторов и SVD малого произведения сохраняют матрицу без потерь,
если ненулевые сингулярные значения не отбрасываются.
Попеременная оптимизация направлений остаётся невыпуклой.\cite{svd,algorithm}

\subsection{Когда допустимо усечение одного SVD}

Если полный квадратичный оператор, включая регуляризацию, разделим:
$\mathcal LX=H XG$ при $H\succ0$, $G\succ0$, то
\eqref{eq:energy} задаёт норму
$\|H^{1/2}(B-B_*)G^{1/2}\|\F^2$.
Обратимая двусторонняя замена координат сохраняет ранг и сводит
задачу к обычной малоранговой аппроксимации Фробениуса.
В общем случае \eqref{eq:operator} является суммой таких операторов,
и это сокращение неприменимо.

Изотропное условие $U_j\T U_j=\eta_j\Id_d$, используемое
в части исходных выводов, особенно ограничительно: при $p<d$
и $\eta_j>0$ оно невозможно по рангу. Даже одинаковая анизотропия
$U_j\T U_j=\eta_jH$ в общем случае не позволяет включить обычный
ridge $\lambda X$ в одно произведение $HXG$.
Поэтому точная спектральная формула в специальном случае не обосновывает
обычное усечение SVD для произвольных ADP-статистик.\cite{svd,unified,archiveaudit}

\subsection{Почему ограничение ранга самой матрицы может мешать}

Поскольку все $m$ сингулярных значений $Y_0$ равны единице,
\begin{equation}\label{eq:floor}
 \min_{\rank(B)\le r}\|B-Y_0\|\F^2=m-r,
 \qquad F_g(B)\ge\lambda(m-r),\quad r<m.
\end{equation}
Кроме того, матрица ранга $r<m$ не содержит полного $m$-мерного
базиса. Дополнение недостающих направлений из $Y_0$ является
использованием предыдущей оценки; эти направления не восстановлены
данными внутри малорангового решения.

Поправка $B=Y_0+\Delta$ избегает такого ограничения на ранг базиса.
Если $Y_0\T\Delta=0$, то $B\T B=\Id+\Delta\T\Delta\succ0$,
и ортонормирование допустимо при любом конечном $\Delta$.
Но фиксация $g_j$, ортонормирование и повторная оценка коэффициентов
не превращают условную задачу \eqref{eq:fixed} в тождественную
профильной задаче \eqref{eq:profile}.\cite{svd,unified}

\subsection{SVD как выбор пространства и результат предобусловливания}

Более осторожная идея из архивных обсуждений состоит в использовании
SVD для выбора пространства, после чего все коэффициенты в нём
решаются совместно. Например, при фиксированном $V\in\R^{d\times q}$
матрица $B=B_0+VK$, $K\in\R^{q\times m}$, определяется линейной
задачей на $qm$ неизвестных. Полное уточнение вне оболочки возможно
через расширенный оператор МНК или методы Крылова.
Условная ортогональность нормальной невязки к выбранной оболочке
не означает, что полный последующий шаг не изменит её коэффициенты:
оператор может связывать оболочку с дополнением.
Для всей предложенной комбинации «SVD~--- малое точное решение~---
полное уточнение» независимого положительного результата полного
обучения в рассмотренных источниках нет.\cite{svd,archivecoarse}

Отдельно проверено диагональное предобусловливание условных
$v$-задач в SVD-решателе. Замена $x=x_0+Sz$, $S\succ0$,
применена к обоим блокам расширенного оператора
$[A;\sqrt\lambda\Id]$. Проверяется исходная нормальная невязка.
При $\lambda=0$ преобразование отключено: оно может изменить
выбор решения минимальной нормы.

В независимой стадии 30 сентября все 144 строки завершились успешно.
На задачах с разными масштабами столбцов медианное отношение времени
к прежней реализации составило $0.225485$, числа проходов по оператору~---
$0.238095$. Максимальное изменение EDR-ошибки полного обучения
составило $2.26\cdot10^{-8}$. Следовательно, подтверждено ускорение
данных условных задач с сохранением результата, а не улучшение
восстановления. Опция \file{precondition_v} по умолчанию выключена.
Устранение вложенных операторных вызовов дало отношение времени
$0.943049$ на обычных фиксированных данных.
Эти результаты не являются проверкой блокового GPU-алгоритма,
который в архиве только предложен.\cite{svdprecond}

\section{Профилирование, повороты и малая модель кривизны}
\label{sec:geometry}

\subsection{Градиент и число плоскостей}

Положим $M_j=U_jY$, $r_j=I_j-M_jg_j$. На гладкой области
профиля без дополнительного штрафа по подпространству
\begin{equation}\label{eq:gradient}
 G=\grad\Phi_\tau(Y)=
 -2(\Id-YY\T)\sum_jN_jU_j\T r_jg_j\T.
\end{equation}
Горизонтальный шаг $Z$ удовлетворяет $Y\T Z=0$. Если
$Z=V\Theta A\T$~--- его сокращённое SVD ранга
$q\le\min(m,d-m)$, то
\begin{equation}\label{eq:exp}
 Y_+=Y(\Id-AA\T)+YA\cos\Theta\,A\T
                         +V\sin\Theta\,A\T.
\end{equation}
Это $q$ поворотов в взаимно ортогональных плоскостях.
Для одного данного касательного шага соответствующие вращения коммутируют.
После пересчёта градиента плоскости обычно меняются, и это свойство
не переносится на всю последовательность итераций.

При $q=1$ разность $Y_+-Y$ имеет ранг не выше единицы,
но $Y_+$ остаётся базисом ранга $m$. Полярная ретракция
\[
 R_Y(Z)=(Y+Z)(\Id+Z\T Z)^{-1/2}
\]
имеет углы $\arctan\theta_a$, тогда как экспонента имеет углы
$\theta_a$. При нескольких различных сингулярных значениях
кривые $R_Y(tZ)$ и $\operatorname{Exp}_Y(tZ)$ в общем случае
не совпадают после одной скалярной замены параметра.\cite{unified,glmcheck}

\subsection{Точная формула углового профиля и её измеренная стоимость}

Для одной плоскости неизменные $m-1$ столбцов дизайна обозначим
$C_j$, а изменяемый~--- $q_j(t)=u_j\cos t+v_j\sin t$.
При $\tau>0$ исключение неизменных коэффициентов даёт
\begin{align}
 \Pi_j&=\Id-C_j(C_j\T C_j+\tau\Id)^{-1}C_j\T,\nonumber\\
 \phi_j(t)&=I_j\T\Pi_jI_j-
 \frac{(q_j(t)\T\Pi_jI_j)^2}{q_j(t)\T\Pi_jq_j(t)+\tau}.
 \label{eq:angle}
\end{align}
Матрица $\Pi_j$ при $\tau>0$ не является ортопроектором.
При $\tau=0$ нужны невырожденность исключаемой системы
и положительный знаменатель; вырожденные точки проверяются отдельно.
После предварительных скалярных произведений оценка суммы
$\sum_jN_j\phi_j(t)$ стоит $O(J)$. Формула вычисляет значение
функции, но не даёт универсального замкнутого решения для оптимального
угла суммы локальных профилей.

Сохранённое измерение для $J=128$, $p=40$, $d=300$, $m=3$,
$\tau=0.2$ и 61 угла сравнивает \eqref{eq:angle} с повторными
пакетными расширенными SVD. На трёх входах отношение времени
с учётом общей проекции равно $0.160$--$0.167$; расхождение кривых
меньше $1.7\cdot10^{-16}$. Измерены медианы семи повторов
после прогрева, один поток BLAS. Это примерно шестикратное ускорение
оценивания уже выбранной кривой, без выбора направления, минимизации
угла или полного обучения. Малая память собственно скалярного
оценивания не означает меньший полный RSS с учётом подготовки.
\cite{theorychecks}

\subsection{Полный дифференциал локальных коэффициентов}

На области, где $A_j=M_j\T M_j+\tau\Id$ обратима, для
$T_j=U_jZ$ дифференцирование локального нормального уравнения даёт
\begin{equation}\label{eq:dg}
 A_j\,dg_j=T_j\T r_j-M_j\T T_jg_j.
\end{equation}
Полный дифференциал расширенной взвешенной невязки имеет вид
\begin{equation}\label{eq:jac}
 d\widetilde r_j=\sqrt{N_j}
 \begin{bmatrix}-T_jg_j-M_jdg_j\\-\sqrt\tau\,dg_j\end{bmatrix},
 \qquad
 \widetilde r_j=\sqrt{N_j}
 \begin{bmatrix}r_j\\-\sqrt\tau\,g_j\end{bmatrix}.
\end{equation}
Член $T_j\T r_j$ нельзя удалять из полного якобиана без изменения
модели кривизны. Множители $\sqrt{N_j}$ необходимы одновременно
в невязке и якобиане.

Выбор $Z=VK$, $K\in\R^{q\times m}$, приводит к малой модели
Гаусса~--- Ньютона на $qm$ переменных. Она полна внутри выбранной
левой оболочки $V$, но не во всём касательном пространстве размерности
$m(d-m)$. Даже $q=m$ в общем случае не устраняет это различие.
Блоковое исключение неактивных коэффициентов через дополнение Шура
может ускорить оценивание модели. Точный ньютоновский шаг в выбранной
карте также не является решением всей задачи на $\mathrm{Gr}(m,d)$.
\cite{unified,structurednote,glmcheck}

\subsection{Ранг и границы гарантий}

При $\tau=0$ обычное обращение $M_j\T M_j$ требует полного
столбцового ранга. Постоянный, но недостаточный ранг позволяет
использовать дифференцирование псевдообратной при соответствующих
условиях, а не формулу с несуществующей обратной матрицей.
Численные пороги ранга дополнительно разбивают область на режимы.
В малом контрпримере профиль равен нулю сколь угодно близко
к вырожденному направлению и единице в самой вырожденной точке.
Поэтому ослабление проверки ранга не является обоснованным ускорением.

Убывание профиля по Армихо не доказывает глобальный минимум
и не обеспечивает сходимости всей последовательности базисов.
Восстановление истинного подпространства требует отдельных
предположения о гладкости, выборе шагов, наблюдаемости и статистической
модели. Эти условия не установлены для общего внешнего ADP-цикла.
\cite{theorychecks,glmcheck}

\section{Реализация и расширение пространства шага: результаты проверок}
\label{sec:negative}

\subsection{Ускорение существующей малой задачи}

Первый переход от теории к вычислениям сохранил малую профильную
задачу, но изменил её реализацию: безопасный QR для $m\le2$
с резервным SVD, ограниченная рабочая память градиента, пакетная
сборка якобиана и повторное использование действий оператора.
Сокращение через QR сохраняет цель лишь при сохранении всей
необходимой информации; если $p<d+1$, QR матрицы $[U_j,I_j]$
само по себе не сокращает число строк.

Сохранённое ускорение фиксированного вызова составило около
$2.52$ раза. В трёх парных полных Spokoiny-обучениях медианное
отношение времени равно $0.941939$, при близких конечных проекторах.
Внутренние вызовы исчерпали пятишаговый бюджет. Некоторые промежуточные
статистики внешних траекторий различались заметно даже при близких
базисах, поэтому результат не следует описывать как побитовое
совпадение всей адаптивной последовательности.
Оптимизированная малая задача вошла в основной решатель;
экспериментальные расширения оценивались уже относительно неё.
\cite{coreoptim}

\subsection{Полный касательный шаг и совместная оптимизация углов}

Полный матрично-свободный GN из исследования \file{glm.tex}
прошёл проверки сопряжения, конечных разностей и сравнение
с плотной расширенной МНК. Однако все 45 измеренных линейных
GN-решений исчерпали лимит LSMR и не прошли сертификат нормальной
невязки. В результате выполнялся прежний малый шаг, а время
фиксированного вызова выросло в $11.63$--$20.21$ раза.
В полных обучениях все 27 внутренних вызовов встретили локальную
потерю ранга; предложенный полный шаг не применялся.
Медианное отношение полного времени составило $1.2143$.
Проверка отвергла данную ограниченную реализацию без предобусловливания;
она не сравнила качество успешно вычисленного полного GN с малым GN.
Отложенные данные не использовались.\cite{glmcheck}

Вариант \file{flat_gn} совместно оптимизировал углы в фиксированной
оболочке сингулярных плоскостей градиента. На 24 фиксированных задачах
медианные отношения конечного профиля и нормы градиента к малому GN
равнялись $0.586$ и $0.606$, времени~--- $2.557$.
В шести парах полного обучения метод был медленнее по медиане
в $1.855$ раза и ухудшил EDR-ошибку в четырёх парах;
наибольший рост составил $0.071155$.
Все внутренние вызовы исчерпали бюджет. В сырых строках полного
обучения отсутствовало отдельное поле конечной ортонормальности;
независимо от этого пробела критерии качества и времени уже не выполнены.
Это прямой пример расхождения улучшения фиксированного профиля
и восстановления.\cite{flat}

\subsection{Дополнение Шура, Ньютон, метрика коэффициентов и операторы CX}

Структурная ветвь из \file{edr_research_note.tex} проверила
исключение неактивных коэффициентов, GN и ньютоновские действия
в малой карте, повторное использование оболочки и метрические шаги.
В основной серии все методы исчерпали 30 шагов: сравнения времени
до общего сертификата не получилось. Медианы времени решателя
составили около $22.1$ мс для малого GN, $45.7$ мс для метрического
метода, $36.4$ мс для Schur GN и $116$--$123$ мс для Ньютона.
Около истинного пространства кандидаты достигали сертификата,
но работали медленнее базового метода.

Положительный частный результат получен в одной дополнительной
модели слабого сигнала: метрический метод сертифицирован на трёх
входах из трёх, базовый~--- на двух. На двух общих сертифицированных
парах ускорение составило $1.6621$ раза; медианные ошибки подпространства
по трём входам равнялись $9.00\cdot10^{-6}$ и $9.22\cdot10^{-4}$.
Это малая калибровка фиксированной задачи, не независимая проверка
полного адаптивного обучения. Schur-ядро вычисления значения
и произведения гессиана на вектор ускорилось примерно в $4.5$ раза
после подготовки; полные обучения от этого быстрее не стали.

Факторизация моментов $U_j=C_jX$ математически точна при согласованных
весах и нормировке. В измерении с $n=400,d=100,J=100,p=8$
отношение объёма хранения факторизованного и плотного представлений
равнялось $0.625$ при опоре из восьми наблюдений и $6.505$
при полной опоре. Во втором случае действия оператора существенно
замедлялись. Экономия зависит от разреженности окон и включает
стоимость построения; одной записи $U=CX$ для вывода об ускорении
недостаточно.\cite{structured}

Последующая проверка метрического спуска на полном Spokoiny-обучении
не прошла отбор: при пяти шагах отношение времени $0.982$
не достигло требуемого $0.909$, а на одном входе ошибка выросла
на $5.666\cdot10^{-4}$ при допуске $10^{-4}$.
На пяти парах фиксированных задач с общим сертификатом
и бюджетом 100 шагов метрический метод был медленнее в $2.040$ раза.
Популяционная OPG-матрица данного генератора оказалась пропорциональной
единичной, поэтому исходная мотивировка через слабое популяционное
направление к нему неприменима.\cite{metric}

\subsection{Совместные блоки вращений и выбор ранга шага}

Проверялись фиксированный, полный, энергетический и зависящий
от стоимости ранги блока. Ни один из четырёх кандидатов не прошёл
порог ускорения $20\%$ на стадии отбора.
Протокол всё же предусматривал отложенное сравнение выбранного
\file{block1}. На десяти парах медианное отношение времени
к малому GN равно $0.99963$, описательный интервал~---
$[0.98914,1.01554]$, медианная разность EDR-ошибок~--- нуль.
Ни один внутренний вызов не достиг стационарности.
При части многоиндексных задач варианты сводились к одной и той же
резервной операции ранга один и потому не различали исходные блочные
гипотезы.

В сохранённых диагностических сравнениях неоднократно меньший профиль
сочетался с большей EDR-ошибкой. Например, при $m=3,d=20,J=64$
профиль $630.4844$ и ошибка $0.54677$ базового метода сменялись
значениями $616.7898$ и $0.56967$ у спектрального блока.
Это наблюдение не оценивает статистический риск вообще, но запрещает
использовать уменьшение внутреннего функционала как единственное
основание выбора ранга шага.\cite{rankk}

\section{Относительный GN: положительная проверка и эффект бюджета}
\label{sec:relative}

\subsection{Предложение на границе ранга}

Диагностика реальных ADP-вызовов показала, что многие проверки ранга
срабатывают из-за фактической локальной вырожденности. При $\tau=0$
и точном равенстве
\[
 \rank(U_j)=\rank(U_jY)<m
\]
пространство столбцов $U_jY$ совпадает с пространством столбцов $U_j$.
Пока равенство рангов сохраняется, локальное значение профиля постоянно.
Это позволяет формировать GN-модель по регулярным центрам.
Численное равенство рангов используется только для построения предложения:
полный градиент, профиль по всем центрам и конечный сертификат
проверяются отдельно.

Принятое изменение масштабирует демпфирование по якобиану регулярных
центров:
\begin{equation}\label{eq:relative}
 \mu=\lambda_{\mathrm{prox}}
         \frac{\|J_{\mathrm{reg}}\|\F^2}{qm}.
\end{equation}
Если все массы или квадрат масштаба отклика меняются в одинаковое число
раз, квадратичная модель и $\mu$ масштабируются согласованно.
Это свойство условной модели; оно не доказывает инвариантность всего
внешнего алгоритма.

Первый вариант intrinsic GN без такой нормировки прошёл отбор,
но провалил независимую проверку качества: агрегированное отношение
EDR-ошибок $0.98546$ сопровождалось одним ухудшением примерно
в $10.069$ раза. Этот отрицательный результат относится к прежнему
варианту и не исключён из истории после появления \eqref{eq:relative}.
\cite{accuracy18,relative}

\subsection{Отбор и независимая проверка relative10}

Для версии \file{relative10} зафиксированы $n=500,m=2$,
$d\in\{5,20,50\}$, десять внутренних шагов с ограничением первых
гладких вызовов пятью шагами, $\tau=0$, один поток BLAS.
Контроль \file{original5} использовал прежний пятишаговый решатель.
На каждой стадии~--- 30 парных задач и 90 обучений трёх методов,
включая фактический R \file{MAVE::MEANMAVE}.

\begin{table}[htbp]
\centering\small
\caption{Relative10 / original5. Интервалы и отношения взяты
из сохранённого парного анализа; меньше единицы означает улучшение.}
\label{tab:relative}
\begin{tabularx}{\linewidth}{@{}Xcc@{}}
\toprule
Стадия, начальные значения генератора & $R_e$ [95\% интервал]
 & Время: медиана / $p_{90}$\\
\midrule
Отбор, 105800--105809 & $0.898376\ [0.844245,0.948903]$ & $0.650580/0.790378$\\
Отложенная, 106800--106809 & $0.923729\ [0.886622,0.956495]$ & $0.653572/0.993297$\\
\bottomrule
\end{tabularx}
\end{table}

Все восемь заранее заданных критериев выполнены на обеих стадиях.
На отложенных данных отношение RSS равно $1.002106$,
максимальное индивидуальное отношение EDR-ошибок~--- $1.036418$.
Подтверждено улучшение при данных бюджетах; интервал не гарантирует
снижение ошибки минимум на $5\%$ в генеральной совокупности.
Эффект сосредоточен при $d=20,50$; при $d=5$ качество почти не меняется.
Во всех сравниваемых ADP-обучениях внешняя остановка~--- $h_{\min}$,
сертифицированных внутренних вызовов нет.
R MEANMAVE сохраняет меньшую ошибку в каждой медиане сценария.
Повторное воспроизведение 120 сохранённых обучений подтвердило
совпадение переноса в публичный код, но не добавило независимых данных.
\cite{relative}

\subsection{Combined20: отделение интеграции от увеличения числа шагов}

Поздняя комбинация добавила прямой угловой поиск по Армихо,
до трёх квадратичных предложений и точное исключение центров с нулевым
оператором, а бюджет увеличила до 20 шагов. Для \file{combined20}
проведены отдельные отбор и независимая проверка; последняя содержит
30 парных задач и 150 обучений пяти методов.

\begin{table}[htbp]
\centering\small
\caption{Независимая проверка combined, начальные значения
108800--108809. Сравнения разделяют состав шага и его бюджет.}
\label{tab:combined}
\begin{tabularx}{\linewidth}{@{}Xcc@{}}
\toprule
Числитель / знаменатель & $R_e$ [95\% интервал]
 & Время: медиана / $p_{90}$\\
\midrule
Combined20 / original5 & $0.925311\ [0.861556,0.983259]$ & $0.860129/1.189895$\\
Combined10 / relative10 & $1.006697\ [0.998577,1.016538]$ & $1.006664/1.038341$\\
Combined20 / combined10 & $0.996743\ [0.974926,1.023747]$ & $1.306993/1.367969$\\
Combined20 / relative10 & $1.003418\ [0.986369,1.028379]$ & $1.275697/1.421531$\\
\bottomrule
\end{tabularx}
\end{table}

Combined20 прошёл исторические критерии относительно original5.
Однако дополнительное улучшение восстановления относительно relative10
не установлено: отношение ошибок близко к единице, интервал допускает
оба знака изменения, а медианное время выше на $27.57\%$.
Удвоение бюджета combined10 также не дало подтверждённого выигрыша
в качестве и повысило медианное время на $30.70\%$.
Прежние эффекты отдельных компонент получены на разных данных
и не могут складываться как независимые проценты.

На отложенной стадии RSS combined20/original5 равен $1.004027$,
максимальная ошибка ортонормальности~--- $4.32\cdot10^{-15}$.
Стационарность сертифицирована лишь для одного из 140 финальных
внутренних вызовов; все внешние остановки~--- $h_{\min}$.
Относительно R MEANMAVE геометрическое отношение EDR-ошибок
combined20 равно $1.296449$ с интервалом $[1.212391,1.393809]$.
Положительный результат здесь~--- улучшение прежнего ADP-контроля,
а не превосходство R по восстановлению.
Параметры \file{max_steps=20}, \file{relative_gn=True},
\file{quadratic_trials=3}, \file{constant_centers=True}
соответствуют текущей сигнатуре Grassmann-решателя.
\cite{combined,currentcode}

\section{Малоранговые поправки, LoRA и проекционный якобиан}
\label{sec:lowrank}

\subsection{Предварительный поиск вариантов после combined20}

Отдельный предварительный поиск проверил девять вариантов:
ортогональную запись полного якобиана, относительное демпфирование
также на гладких задачах, диагональные плоскости кривизны,
интерполяцию длины GN-шага, три положительных локальных ridge
и два коротких полных шага Крылова.
В четырёх сериях записано 80 обучений, но использованы одни и те же
четыре сценария с начальным значением 90800.
Ни один кандидат не показал совместного выигрыша качества и времени.

Например, отношения EDR/времени к combined20 составили
$1.0000/1.032$ для \file{stable_jacobian}, $1.1713/1.049$
для диагонального варианта, $1.0282/0.838$ для ridge $10^{-6}$
и $1.0304/1.444$ для \file{krylov20}.
Положительный ridge сделал профиль гладким ценой изменения
оценивающего функционала и ухудшения восстановления.
Алгебраическая перезапись якобиана сама по себе не дала измеримого
выигрыша точности: приведённые ошибки относительно расчёта повышенной
точности близки, около $3.81\cdot10^{-9}$ и $3.84\cdot10^{-9}$.
Новые данные отбора и независимой проверки в этой серии не использовались.
\cite{screen22}

\subsection{Что даёт фиксирование коэффициентов}

При текущих профильных $g_j$ и $r_j$ поздний SVD-подход решает
суррогатную задачу
\begin{equation}\label{eq:lowrank}
 q(\Delta)=\sum_jN_j\|r_j-U_j\Delta g_j\|^2+\mu\|\Delta\|\F^2,
 \quad Y\T\Delta=0,\quad\rank\Delta\le k.
\end{equation}
Проверены жадная SVD-поправка, совместное решение всех коэффициентов
в сингулярной оболочке градиента, двухфакторная LoRA-подобная
параметризация и обогащение оболочки полной нормальной невязкой.

Для $\tau=0$, $\mu\ge0$ и точной арифметики справедливо
\begin{equation}\label{eq:surrogate}
 \Phi_0\bigl(\polar(Y+\Delta)\bigr)
 \le\sum_jN_j\|r_j-U_j\Delta g_j\|^2\le q(\Delta).
\end{equation}
Действительно, $Y+\Delta=Y_+R$ с обратимой $R$;
коэффициенты $Rg_j$ воспроизводят ту же невязку на $Y_+$,
а последующее профилирование может только уменьшить её.
Это достаточная связь суррогата с предложением спуска, если удалось
уменьшить $q$ относительно $q(0)$. Она не доказывает глобальный
минимум задачи ранга $k$ или улучшение EDR. При $\tau>0$
перенос коэффициентов изменяет ridge-штраф.
В реализации дополнительно проверяется исходный профиль с его
численными порогами ранга.\cite{lowrankmath}

\subsection{Пилот малоранговых и LoRA-подобных методов}

В первом пилоте использованы семь сценариев и одно начальное значение
генератора 90800: всего 77 обучений одиннадцати методов.
Это семь общих наборов данных, а не 77 независимых наблюдений.
Для \file{lora_k10} получены $R_e=0.990$ и отношения времени
$1.017/1.164$ (медиана/$p_{90}$); при $d=50$ ошибка выше на $4.7\%$.
Для \file{svd_k10} время снижено до $0.660$ от контроля,
но $R_e=1.014$, а в коррелированном сценарии ошибка выше на $15.1\%$.
Внутренний LSMR LoRA-вариантов часто исчерпывал лимит трёх итераций.

Ранг $k=\min(m,d-m)$ в новых полноразмерных оболочках
не соответствует ограничению $r<m$ буквального старого жадного SVD.
Это не сравнение двух алгоритмов при одном и том же ранге допустимой
поправки. Независимой отложенной проверки LoRA в этих материалах нет.
\cite{lowrank}

\subsection{Проекционная модель: точный градиент, другая кривизна}

При $\tau=0$ и полном столбцовом ранге $M_j$ дифференциал
профильной невязки раскладывается на ортогональные компоненты:
\begin{equation}\label{eq:projected}
 dr_j=-(\Id-M_jM_j^\dagger)T_jg_j
      -(M_j^\dagger)\T T_j\T r_j.
\end{equation}
Первая лежит в дополнении к столбцам $M_j$, вторая~--- в их оболочке.
Поскольку $M_j\T r_j=0$, сохранение только первой компоненты
оставляет тот же градиент. Но для взвешенных якобианов
\[
 J_{\mathrm{full}}\T J_{\mathrm{full}}
 -J_\perp\T J_\perp=J_\parallel\T J_\parallel\succeq0.
\]
Следовательно, это приближённая модель Гаусса~--- Ньютона
с отброшенной частью кривизны, а не более точная запись полного якобиана.
\cite{lowrankmath}

\subsection{Замороженный отбор и контроль вычислительного бюджета}

После благоприятного пилота с одним начальным значением был зафиксирован отбор
на 70 новых парных задачах: семь сценариев и десять значений
111800--111809. Четыре метода дали 280 успешных обучений:
combined20, тот же решатель с десятью шагами (\file{budget10}),
\file{projected10} и R MEANMAVE. Интервалы рассчитаны по 10000
ресэмплированиям блоков начальных значений с сохранением семи сценариев
внутри блока.

\begin{table}[htbp]
\centering\small
\caption{Отбор проекционной модели. Первые две строки используют
один контроль baseline20; последняя отделяет изменение кривизны от бюджета.}
\label{tab:projected}
\begin{tabularx}{\linewidth}{@{}Xcc@{}}
\toprule
Сравнение & $R_e$ [95\% интервал] & Время: медиана / $p_{90}$\\
\midrule
Budget10 / baseline20 & $0.998953\ [0.976752,1.024110]$ & $0.742179/0.941168$\\
Projected10 / baseline20 & $1.000646\ [0.981306,1.022270]$ & $0.7366/0.8849$\\
Projected10 / budget10 & $1.001694\ [0.988592,1.016779]$ & $0.990/1.029$\\
\bottomrule
\end{tabularx}
\end{table}

Требовались $R_e\le0.98$, верхняя граница интервала ниже единицы,
рост медианной ошибки в каждом сценарии не более $0.005$
и индивидуальный рост не более $0.05$.
Обе десятишаговые версии не прошли эти критерии качества.
Для projected10 наибольший рост медианы равен $0.0110$,
индивидуальный рост~--- $0.1238$; критерии времени, RSS
и ортонормальности выполнены.
Почти такое же ускорение budget10 показывает, что общий выигрыш
во времени нельзя приписать проекционной кривизне.
Её добавочного улучшения восстановления относительно budget10 не выявлено.

Отложенные значения 112800--112809 не использовались,
изменение не принято. LoRA-метода в этом отборе не было:
отрицательный результат projected10 не является проверкой LoRA.
Все три серии вместе содержат 385 обучений на 77 различных
парах «сценарий, начальное значение»; пилотные данные переиспользованы.
\cite{lowrank}

\section{Анизотропия и центры: проверка статистических изменений}
\label{sec:centers}

\subsection{Идеи локализатора и критериев отбора}

Отдельная рукопись предлагает ограничивать снизу собственные значения
локализатора, отбирать кандидатов дешёвым поиском в текущей проекции,
а центры выбирать по LogDet-критерию с массовым покрытием.
Ограничение спектра, множитель плотности и сокращение центров
меняют статистическую процедуру. Точный предварительный фильтр
при неизменной геометрии может сохранять опору ядра, но сокращение
набора центров вообще не сохраняет исходный функционал.

Для фиксированных векторов $u_j$ и $\varepsilon>0$
\[
 f(S)=\log\det\left(\varepsilon\Id+
                  \sum_{j\in S}u_ju_j\T\right)
             -\log\det(\varepsilon\Id)
\]
монотонен и субмодулярен. При ограничении мощности множества
обычная жадная процедура имеет соответствующую гарантию приближения
\emph{этого критерия}. Строгого роста на каждом добавлении нет,
например при $u_j=0$. Если векторы или штрафы перекрытия изменяются
в зависимости от уже выбранного множества, свойство надо доказывать
заново. Оно не даёт гарантии EDR-восстановления.

Аналогично точное сохранение суммарной массы квадратурой разбиения
единства не означает сохранения всех нормальных операторов:
последние зависят от того, как локальные матрицы меняются между центрами.
Эти ограничения сформулированы в самой поздней ревизии рукописи,
но её рекомендации о готовности к внедрению сильнее имеющихся проверок.
\cite{centerstex,centersprotocol}

\subsection{Числа исходной рукописи и пробелы происхождения}

В таблицах рукописи заявлены $n=1000$, $J=500$, $p=40$,
пять значений 101--105, общие данные, инициализация и направления.
При увеличении спектрального порога от нуля до $0.50$
приведены изменения медианной ошибки:
\[
 \begin{array}{ll}
 \text{additive\_d30:}&0.0280\to0.0052,\\
 \text{multiplicative\_d30:}&0.0975\to0.0304,\\
 \text{aniso\_curv\_d40:}&0.1819\to0.0161,\\
 \text{additive\_d50:}&0.1773\to0.0292.
 \end{array}
\]
Это \emph{числа исходного текста}. Ссылок на сырые прогоны,
их конфигурации и точное определение ошибки в экспериментальном
разделе нет. Поэтому они не включаются в свод подтверждённых
улучшений.

Даже внутри таблиц формулировка о монотонном улучшении по порогу
не выполняется: в multiplicative\_d30 медиана при пороге $0.20$
равна $0.1688$, что хуже $0.0975$ при нулевом пороге.
Для одинаково подписанного SF-ADP с порогом $0.25$
приведены разные медианы: aniso\_curv\_d40~--- $0.0424$
и $0.0399$, additive\_d50~--- $0.0427$ и $0.0525$
в разных таблицах. Без первичных данных причины расхождений
установить нельзя.\cite{centerstex}

\subsection{Последующая воспроизводимая проверка}

Существует отдельное исследование по мотивам этой рукописи.
Оно прямо указывает, что исходные сырые данные недоступны,
и не является их воспроизведением. Новый протокол использует
$n=500$, $m=2$, $p=10$, четыре сценария, три значения
92800--92802, десять вариантов и штатный LSMR.
Все 120 обучений завершились, парность входов и инициализации
проверена по хешам. Ни в одном обучении не сошлись все внутренние
сертификаты; остановки по $h_{\min}$ не интерпретируются как сходимость.

\begin{table}[htbp]
\centering\small
\caption{Адаптивные центры: новый отбор, 12 пар на вариант.
Разность EDR~--- кандидат минус контроль, время включает отбор центров.}
\label{tab:centers}
\begin{tabularx}{\linewidth}{@{}Xccc@{}}
\toprule
Вариант / контроль & Медиана $\Delta e$ & Время & $e\le0.1$\\
\midrule
Floor25 / baseline & $0.0000$ & $0.994$ & $1/12$\\
Floor50 / baseline & $+0.0015$ & $0.989$ & $2/12$\\
Поиск проекции / floor25 & $\approx0$ & $1.682$ & $1/12$\\
LogDet / baseline & $+0.0464$ & $2.481$ & $0/12$\\
LogDet + покрытие / baseline & $+0.0397$ & $2.487$ & $0/12$\\
LogDet + покрытие + квадратура / baseline & $+0.0242$ & $2.814$ & $0/12$\\
\bottomrule
\end{tabularx}
\end{table}

Floor25 не прошёл заранее заданный критерий улучшения:
медианная разность не достигла $-0.01$, одна пара ухудшилась
больше чем на $0.05$. Выигрыш в additive30 не переносится
на другие сценарии. Floor50 был дополнительным вариантом без
заранее объявленного критерия допуска. Отбор центров ухудшил
совокупное качество и время; сокращение числа центров также
не уменьшило RSS процесса. Отложенные значения 92900--92904
не запускались.

Различия $n$, $p$, генераторов и настроек запрещают трактовать
эту серию как опровержение чисел на тех же данных.
Она показывает, что заявленный широкий положительный эффект
не подтверждён новым зафиксированным протоколом.
Даже совпадение опор и статистик до машинной точности не гарантирует
одинаковых адаптивных траекторий: отдельный сохранённый аудит
обнаружил усиление различий после первого вызова решателя,
не установив единственную причину.\cite{centersreport,centersprotocol}

\section{Тензорные поля, статистические метрики и открытые гипотезы}
\label{sec:open}

\subsection{Аппроксимация оператора и другие правила выбора шага}

Предложено приближать нормальный оператор \eqref{eq:operator}
одним или несколькими разделимыми слагаемыми. Для положительно
определённого приближения $\mathcal L_0$ условие
\[
 \rho=\|\mathcal L_0^{-1/2}(\mathcal L-\mathcal L_0)
                     \mathcal L_0^{-1/2}\|_2<1
\]
даёт оценку обусловленности предобусловленной системы
$(1+\rho)/(1-\rho)$. Требуется именно оценка операторной нормы;
нескольких случайных проб для её доказательства недостаточно.
Кроме того, SVD не обеспечивает положительности факторов приближения.
Оператор $\mathcal L_0$ можно использовать для предобусловливания
исходной системы. Замена им самого целевого оператора определяет
уже другую задачу. Общий выигрыш полного обучения для этой схемы
не подтверждён.
\cite{unified}

В архиве также обсуждаются оптимизация факторов при фиксированном
ранге, ядерный штраф и выбор направления в нормах Шаттена или Кая Фана.
Факторные координаты позволяют сократить число переменных,
но задача остаётся невыпуклой. Ядерный штраф меняет оцениватель.
Выбранная норма определяет двойственную задачу линейного спуска.
Её решение не обязано иметь ранг один. В частности,
в шаре спектральной нормы при градиенте $\Id_2$ оптимум линейного
спуска может иметь ранг два. Для этих общих семейств в корпусе
нет самостоятельного положительного полного сравнения.
\cite{svd,glmcheck}

\subsection{CP/Tucker и пространственно меняющееся подпространство}

Вместо одного глобального пространства предлагается оценивать
поле $[Y(x)]$. Разложение CP/Tucker может сжимать его линейные
коэффициенты либо касательные поправки к общему базису $Y_0$.
Во втором случае условие $Y_0\T Z(x)=0$ позволяет применять
ретракцию отдельно в каждой точке, сохраняя полный локальный ранг.

Малый тензорный ранг $Z(x)$ не означает такого же ранга тензора
конечных ортонормированных базисов: ретракция нелинейна.
Если все базисы лежат в одной оболочке размерности $R=m$,
их $m$-мерные подпространства совпадают. Для нетривиального
вращающегося поля в общей оболочке необходимо $R>m$.
Произвольное гладкое поле подпространств также не обязано быть
касательным полем некоторого многообразия предикторов:
интегрируемость требует отдельных условий.
Полного EDR-сравнения с независимым подтверждением для предложенных
Tucker-полей в рассмотренном корпусе нет.\cite{tucker,unified,structurednote}

\subsection{OPG, MAVE, ковариационные степени и GLS}

Пусть локальный момент задаётся оператором направлений $S_j$:
$U_j=S_j\Sigma_j$, $I_j=S_jc_j$.
Тогда невязка измеряется в метрике
\[
 \|S_j(c_j-\Sigma_jYg)\|^2
  =(c_j-\Sigma_jYg)\T S_j\T S_j(c_j-\Sigma_jYg).
\]
Она не совпадает автоматически с критерием локальной регрессии
MAVE или с евклидовой ошибкой оценки градиента.
В моделях интерполяции через $\Sigma_j^\kappa$ изменение $\kappa$
обычно меняет оцениватель. Степень вырожденной ковариации
не восстанавливает информацию из её ядра.

Для GLS нужна ковариация шума именно используемых моментов.
При условно независимых ошибках наблюдений с дисперсией $\sigma^2$
ковариация нормированного локального момента содержит
\[
 \frac{\sigma^2}{N_j^2}
 \sum_i w_{ij}^{\,2}(x_i-\bar x_j)(x_i-\bar x_j)\T,
\]
а не ковариацию с первыми степенями $w_{ij}$.
Перекрывающиеся окна дополнительно дают межцентровые зависимости.
Поэтому названия «GLS» и «информационная метрика» требуют
уточнённой вероятностной модели.
Гомотопии OPG--MAVE--ADP и такие метрики остаются статистическими
гипотезами; независимой успешной проверки их полного адаптивного
варианта в данном корпусе не найдено.\cite{unified,glm,structurednote}

\subsection{Наблюдаемость и сходимость}

При бесшумном профиле, $\tau=0$ и полном локальном ранге
линейная наблюдаемость касательного направления $Z$ определяется формой
\[
 \mathcal I(Z)=\sum_jN_j
 \|(\Id-M_jM_j^\dagger)U_jZg_j^*\|^2.
\]
Для её невырожденности необходим размерностный запас
$J(p-m)\ge m(d-m)$, но его недостаточно: локальные коэффициенты
и операторы должны возбуждать все направления.
Случай $J=1,m=1,p<d-1$ уже исключает инъективность по размерности,
независимо от гауссовости направлений.
Численные проверки формул не устанавливают положительной нижней
границы этой формы для адаптивной случайной последовательности.
\cite{structurednote}

\section{Свод выводов}
\label{sec:conclusions}

Математически обоснованы операторная форма условного шага,
точный масштаб выбранного атома, условия двустороннего спектрального
решения, профилирование локальных коэффициентов, разложение касательного
шага на повороты и малые модели Шура. Эти результаты применимы в указанных областях ранга
и регуляризации; общую теорию восстановления они не заменяют.

Подтверждённое ускорение без существенного изменения условного результата
получено для диагонального предобусловливания SVD-подзадач.
Для некоторых точных преобразований измерено ускорение отдельных
операций; улучшение полного обучения из него не следует.
Подтверждённое улучшение качества
и времени полного ADP относительно прежнего контроля получено
для relative10 на независимых двухиндексных задачах.
Combined20 также лучше старого контроля, но его добавочный эффект
относительно relative10 не установлен.

Расширение числа переменных шага само по себе не оправдалось:
полный GN, flat GN, метрические и блочные варианты не прошли
соответствующие критерии стоимости или качества. Projected10
не улучшил восстановление относительно контроля с тем же бюджетом.
LoRA-поправки и тензорные поля сохраняют исследовательский статус.
Рекомендации о широком выигрыше адаптивных центров не подтверждены
последующим зафиксированным сравнением.

Главное ограничение корпуса~--- разрыв между внутренней аппроксимацией
и внешним статистическим результатом. Снижение профиля, достижение
градиентного допуска, малый размер линейной задачи и ускорение отдельной
операции отвечают на разные вопросы. В рассмотренных экспериментах
ни один из этих признаков по отдельности не служит достаточным
основанием для вывода об улучшении EDR-восстановления.

\appendix
\section{Карта происхождения идей}\label{sec:corpus}

\small
\begin{longtable}{@{}>{\raggedright\arraybackslash}p{0.31\linewidth}>{\raggedright\arraybackslash}p{0.64\linewidth}@{}}
\toprule
Материал в \file{SVD/} & Содержание и роль\\
\midrule
\endhead
\file{archive/SVD.tex}, \file{archive/SVD_form.tex}
 & Исходные условные малоранговые схемы; жадные атомы, блоки,
 спектральные частные случаи. Поздние ограничения собраны в
 \file{archive/problem.md} и \file{SVD_solver.tex}.\\
\file{archive/chat_1.md}, \file{archive/chat_3.md}
 & SVD как выбор оболочки с последующим полным уточнением;
 совместные блоки и GPU. \file{archive/chat_2.md} пуст.\\
\file{SVD_solver.tex}, \file{SVD_algorithm.tex}, \file{SVD_cur.tex}
 & Систематическая теория условного решателя, последовательность
 алгоритма и историческое описание реализации; не три независимых эксперимента.\\
\file{grassmann_rank_one.tex}, \file{cladue.md}
 & Профильный поворот ранга один; предложения по геометрии,
 нормам и выбору нескольких плоскостей.\\
\file{argon.md}, \file{kimi.tex}, \file{glm.tex}
 & Объединяющие постановки, плоские повороты, полная кривизна,
 ковариационные метрики и открытые статистические расширения.\\
\file{EDR_unified_theory.tex}
 & Исправленная общая теория; разграничение рангов, метрик
 и областей применимости. Проверки хранятся отдельно.\\
\file{edr_research_note.tex}
 & Schur-исключение, производные, оператор $CX$, повторное
 использование оболочки; источники структурной ветви экспериментов.\\
\file{Tucker.md}
 & Пространственное поле, CP/Tucker, общая оболочка и локальные ядра.\\
\file{adaptive_centers_anisotropic_adp_edr_report.tex}
 & Спектральный порог, проекционный фильтр, LogDet, покрытие,
 квадратура; теоретические положения и таблицы без первичных прогонов.\\
\file{multiindex.tex}
 & Базовая модель и ADP/AO-процедура; контекст последующих изменений.\\
\file{GRASSMANN_IMPROVEMENT_HISTORY.md}
 & Вторичный путеводитель по экспериментам. Численные выводы
 настоящего документа сверены с первичными отчётами и избранными сырыми строками.\\
\bottomrule
\end{longtable}
\normalsize

\section{Проверяемость и отсутствующие свидетельства}\label{sec:audit}

В \file{SVD/documentation/math_checks.json} сохранены 24 показателя
алгебраических проверок условного решателя; они имеют разный смысл
и не все являются погрешностями. В
\file{SVD/theory_2026_10_01/formula_results.json} сохранены 203 проверки
в 35 группах, включая экспоненту, профиль, производные, метрики
и контрпримеры. Их скрипты выполнены повторно для этой систематизации
в отдельном каталоге; исходные результаты не перезаписывались.
Проверки пройдены. Они подтверждают малые численные тождества,
но не производительность и не статистическое восстановление.

В \file{kimi.tex} приведено «сохранённое сравнение 350 прогонов»:
ошибка $0.118$ и время $24.8$ с для SVD ранга один против
$0.0464$ и $37.1$ с для LSMR. Не указаны путь к данным,
единица числа прогонов, точная метрика и границы таймера.
Эти значения не используются для вывода о сравнительной эффективности.
Сравнение с 200 случайными матрицами, также названное там
подтверждением глобального минимума, проверяет только эти 200
кандидатов; глобальность следует из условий спектрального вывода,
а не из такого теста.

В \file{edr_research_note.tex} упомянут файл
\file{edr_note_checks.py}, отсутствующий в данной рабочей копии.
Его заявленные погрешности не отождествляются с другими сохранёнными
наборами проверок. Поздние структурные эксперименты имеют собственные
отчёты и тесты. Исходные таблицы адаптивных центров по пяти реализациям
также не получают подтверждения от нового эксперимента с иными данными.

Подробная запись текущей проверки, включая источники таблиц,
пересчёт ключевых парных агрегатов, хеши файлов и сборку документа,
находится в \file{SVD/synthesis_2026_10_24/VERIFICATION.md}
и \file{evidence_audit.json} рядом с ней.

\begin{thebibliography}{99}
\small
\bibitem{svd} Условный SVD-решатель: постановка, доказательства,
ограничения и направления развития.
\repo{SVD/SVD_solver.tex}.
\bibitem{algorithm} Последовательность алгоритма.
\repo{SVD/SVD_algorithm.tex}.
\bibitem{archiveaudit} Критический разбор исходной малоранговой схемы.
\repo{SVD/archive/problem.md}.
\bibitem{archivecoarse} Оболочка SVD и полное уточнение.
\repo{SVD/archive/chat_1.md}.
\bibitem{unified} Общая теория разделимых моделей, поворотов и метрик.
\repo{SVD/EDR_unified_theory.tex}.
\bibitem{theorychecks} Независимые формульные проверки и узкое
измерение угловой кривой.
\repo{SVD/theory_2026_10_01/REPORT.md};
\repo{SVD/theory_2026_10_01/formula_results.json};
\repo{SVD/theory_2026_10_01/angle_results.json}.
\bibitem{svdprecond} Предобусловливание SVD-подзадач, 30 сентября.
\repo{experiments/svd_improvements_2026-09-30/REPORT.md};
\repo{experiments/svd_improvements_2026-09-30/validation_summary.json}.
\bibitem{glm} Предложения по плоским поворотам и расширению шага.
\repo{SVD/glm.tex}.
\bibitem{glmcheck} Проверка полного матрично-свободного GN, 2 октября.
\repo{experiments/glm_grassmann_2026_10_02/REPORT.md};
\repo{experiments/glm_grassmann_2026_10_02/PROOF.md}.
\bibitem{coreoptim} Оптимизация малой задачи и её перенос в основной решатель.
\repo{experiments/grassman_optim_2026_10_01/REPORT.md};
\repo{experiments/grassman_promotion_2026_10_03/REPORT.md}.
\bibitem{flat} Совместная оптимизация углов: фиксированные задачи
и полное обучение.
\repo{experiments/grassman_flat_2026_10_02/REPORT.md}.
\bibitem{structurednote} Исследовательская записка о Schur, метриках и CX.
\repo{SVD/edr_research_note.tex}.
\bibitem{structured} Структурная ветвь: протокол и результаты.
\repo{experiments/edr_structured_2026_10_15/REPORT.md}.
\bibitem{metric} Метрический спуск на полном Spokoiny-обучении.
\repo{experiments/metric_spokoiny_2026_10_15/REPORT.md}.
\bibitem{rankk} Блочные повороты и выбор ранга шага.
\repo{experiments/grassman_rank_k/REPORT.md}.
\bibitem{accuracy18} Диагностика ранга и последовательность вариантов.
\repo{experiments/grassman_accuracy_2026_10_18/REPORT.md};
\repo{experiments/grassman_accuracy_2026_10_18/INTRINSIC_REPORT.md}.
\bibitem{relative} Relative GN: отбор, независимая проверка и перенос.
\repo{experiments/grassman_accuracy_2026_10_18/RELATIVE_GN_REPORT.md};
\repo{benchmark_outputs/r_sdr/grassman-relative-heldout-2026-10-22/analysis_summary.json}.
\bibitem{combined} Combined20: контролируемая проверка интеграции и бюджета.
\repo{experiments/grassman_accuracy_2026_10_18/COMBINED_REPORT.md};
\repo{benchmark_outputs/r_sdr/grassman-combined-heldout-2026-10-22/combined_analysis.json}.
\bibitem{lowrankmath} Поправки и проекционный якобиан: математическая модель.
\repo{experiments/grassman_accuracy_2026_10_22/LOWRANK_MATH.md}.
\bibitem{screen22} Предварительное сравнение девяти вариантов.
\repo{experiments/grassman_accuracy_2026_10_22/REPORT.md};
\repo{experiments/grassman_accuracy_2026_10_22/MATH.md}.
\bibitem{lowrank} Малоранговый пилот и неудачный отбор projected10.
\repo{experiments/grassman_accuracy_2026_10_22/LOWRANK_REPORT.md};
\repo{benchmark_outputs/r_sdr/grassman-lowrank22-development/paired_analysis.json}.
\bibitem{metricdef} Определение EDR и протокол сравнения с R.
\repo{experiments/r_sdr/design.py}; \repo{experiments/r_sdr/README.md}.
\bibitem{centerstex} Рукопись об анизотропии и центрах.
\repo{SVD/adaptive_centers_anisotropic_adp_edr_report.tex}.
\bibitem{centersprotocol} Новый протокол проверки идей адаптивных центров.
\repo{experiments/adaptive_centers/PROTOCOL.md}.
\bibitem{centersreport} Адаптивные центры: зафиксированный отбор 120 обучений.
\repo{experiments/adaptive_centers/REPORT.md};
\repo{experiments/adaptive_centers/results/fullfit_dev.jsonl}.
\bibitem{tucker} Пространственные CP/Tucker-модели.
\repo{SVD/Tucker.md}.
\bibitem{currentcode} Проверенные при систематизации публичные сигнатуры:
\repo{ADP/solver/grassman.py}, \repo{ADP/solver/SVD.py}.
Полное описание действующей процедуры:
\repo{tex/ADP_Grassmann_current.tex}.
\end{thebibliography}
\end{document}
